% ==================== SECCIÓN COMPLETA: MOC ====================
% Incluir en el documento principal después de la sección de Inharmonicidad

\selectlanguage{spanish}
\subsection{MOC (Modal Octave Compression)}

\subsubsection{Definición y Contexto Físico}

La Compresión Modal de Octava (MOC) es una métrica que cuantifica el comportamiento no lineal del instrumento entre la primera y segunda octava. Fue introducida por Benade y otros investigadores para caracterizar cómo se comportan los modos acústicos en instrumentos de viento.

\subsubsection{Derivación Completa}

Partimos de la relación entre frecuencia y longitud efectiva para un tubo abierto en ambos extremos:

\begin{equation}
f_n = \frac{n c_0}{2 L_{\text{eff}}}
\label{eq:frecuencia_tubo}
\end{equation}

donde $n$ es el número del modo, $c_0$ la velocidad del sonido, y $L_{\text{eff}}$ la longitud efectiva del tubo.

Para la primera octava ($n=1$):
\begin{equation}
\freqplay = \frac{c_0}{2 L_{\text{eff},I}}
\label{eq:freq_play_I}
\end{equation}

Para la segunda octava ($n=2$):
\begin{equation}
2\freqplay = \frac{c_0}{L_{\text{eff},II}}
\label{eq:freq_play_II}
\end{equation}

Sin embargo, las frecuencias medidas $\freqzero$ y $\freqone$ no coinciden exactamente con estas frecuencias teóricas debido a efectos de extremos y agujeros.

Definimos las desviaciones:
\begin{align}
\Delta L_I &= L_{\text{eff},I} - L_{\text{real}} = \frac{c_0}{2}\left(\frac{1}{\freqplay} - \frac{1}{\freqzero}\right) \\
\Delta L_{II} &= L_{\text{eff},II} - L_{\text{real}} = c_0\left(\frac{1}{2\freqplay} - \frac{1}{\freqone}\right)
\end{align}

La diferencia entre estas desviaciones es:
\begin{equation}
\Delta\Delta L = \Delta L_{II} - \Delta L_I = c_0\left[\left(\frac{1}{2\freqplay} - \frac{1}{\freqone}\right) - \left(\frac{1}{\freqplay} - \frac{1}{\freqzero}\right)\right]
\end{equation}

Simplificando:
\begin{align}
\Delta\Delta L &= c_0\left[\frac{1}{2\freqplay} - \frac{1}{\freqone} - \frac{1}{\freqplay} + \frac{1}{\freqzero}\right] \\
&= c_0\left[\frac{1}{\freqzero} - \frac{1}{\freqone} - \frac{1}{2\freqplay}\right] \\
&= c_0\left[\left(\frac{1}{\freqzero} - \frac{1}{\freqplay}\right) - \left(\frac{1}{\freqone} - \frac{1}{2\freqplay}\right)\right]
\end{align}

El MOC se define como la razón entre estas diferencias de períodos:
\begin{equation}
\text{MOC} = \frac{\frac{1}{\freqone} - \frac{1}{2\freqplay}}{\frac{1}{\freqzero} - \frac{1}{\freqplay}}
\label{eq:moc_def}
\end{equation}

\subsubsection{Interpretación Física}

\begin{itemize}
    \item \textbf{MOC = 1.0}: Compresión ideal, comportamiento lineal perfecto
    \item \textbf{MOC > 1.0}: Sobre-compresión, la segunda octava está más comprimida
    \item \textbf{MOC < 1.0}: Sub-compresión, la segunda octava está menos comprimida
\end{itemize}

El MOC está relacionado con:
\begin{itemize}
    \item Efectos de extremos abiertos
    \item Comportamiento de agujeros laterales
    \item Acoplamiento entre modos acústicos
    \item Efectos no lineales del flujo
\end{itemize}

\selectlanguage{french}
\subsection{MOC (Compression Modale d'Octave)}

\subsubsection{Définition et Contexte Physique}

La Compression Modale d'Octave (MOC) est une métrique qui quantifie le comportement non linéaire de l'instrument entre la première et la deuxième octave. Elle a été introduite par Benade et d'autres chercheurs pour caractériser comment se comportent les modes acoustiques dans les instruments à vent.

% ==================== SECCIÓN COMPLETA: B_I ====================

\selectlanguage{spanish}
\subsection{B$_I$ (Bias de la Primera Octava)}

\subsubsection{Definición}

B$_I$ mide la desviación de la frecuencia de resonancia medida $\freqzero$ respecto a la frecuencia teórica de ejecución $\freqplay$:

\begin{equation}
B_I = 1200 \log_2\left(\frac{\freqplay}{\freqzero}\right) \quad \text{[cents]}
\label{eq:bi_def}
\end{equation}

\subsubsection{Derivación desde la Teoría de Longitud Efectiva}

Para un tubo abierto en ambos extremos, la frecuencia fundamental es:
\begin{equation}
\freqplay = \frac{c_0}{2 L_{\text{eff}}}
\end{equation}

donde $L_{\text{eff}}$ incluye correcciones de extremos:
\begin{equation}
L_{\text{eff}} = L_{\text{geom}} + \Delta L
\end{equation}

La corrección $\Delta L$ se debe a:
\begin{itemize}
    \item Efectos de extremo abierto (end correction)
    \item Efectos de agujeros abiertos
    \item Acoplamiento con el exterior
\end{itemize}

Si la frecuencia medida $\freqzero$ difiere de $\freqplay$, esto indica una longitud efectiva diferente:
\begin{equation}
\freqzero = \frac{c_0}{2 L_{\text{eff},0}}
\end{equation}

La relación entre frecuencias es:
\begin{equation}
\frac{\freqplay}{\freqzero} = \frac{L_{\text{eff},0}}{L_{\text{eff}}}
\end{equation}

Tomando logaritmo en base 2 y multiplicando por 1200:
\begin{align}
B_I &= 1200 \log_2\left(\frac{\freqplay}{\freqzero}\right) \\
&= 1200 \log_2\left(\frac{L_{\text{eff},0}}{L_{\text{eff}}}\right)
\end{align}

\subsubsection{Interpretación}

\begin{itemize}
    \item \textbf{$B_I = 0$ cents}: Afinación perfecta, $\freqzero = \freqplay$
    \item \textbf{$B_I > 0$ cents}: Nota suena más aguda de lo esperado, $\freqzero < \freqplay$
    \item \textbf{$B_I < 0$ cents}: Nota suena más grave de lo esperado, $\freqzero > \freqplay$
\end{itemize}

\selectlanguage{french}
\subsection{B$_I$ (Biais de la Première Octave)}

% ==================== SECCIÓN COMPLETA: ESPE ====================

\selectlanguage{spanish}
\subsection{ESPE (Effective Sound Path Extension)}

\subsubsection{Definición}

ESPE cuantifica la diferencia en la extensión efectiva del camino del sonido entre la primera y segunda octava:

\begin{equation}
\text{ESPE} = 1200 \log_2\left(\frac{L_{\text{eff},I}}{L_{\text{eff},I} + \Delta\Delta L}\right) \quad \text{[cents]}
\label{eq:espe_def}
\end{equation}

\subsubsection{Derivación Completa}

De la ecuación~\eqref{eq:freq_play_I}, la longitud efectiva para la primera octava es:
\begin{equation}
L_{\text{eff},I} = \frac{c_0}{2 \freqplay}
\end{equation}

La desviación de longitud para la primera octava es:
\begin{equation}
\Delta L_I = \frac{c_0}{2}\left(\frac{1}{\freqplay} - \frac{1}{\freqzero}\right)
\end{equation}

Para la segunda octava:
\begin{equation}
\Delta L_{II} = c_0\left(\frac{1}{2\freqplay} - \frac{1}{\freqone}\right)
\end{equation}

La diferencia entre desviaciones es:
\begin{align}
\Delta\Delta L &= \Delta L_{II} - \Delta L_I \\
&= c_0\left[\left(\frac{1}{2\freqplay} - \frac{1}{\freqone}\right) - \left(\frac{1}{\freqplay} - \frac{1}{\freqzero}\right)\right] \\
&= c_0\left[\frac{1}{\freqzero} - \frac{1}{\freqone} - \frac{1}{2\freqplay}\right]
\end{align}

La longitud efectiva corregida para la segunda octava es:
\begin{equation}
L_{\text{eff},II} = L_{\text{eff},I} + \Delta\Delta L
\end{equation}

ESPE mide el cambio relativo:
\begin{align}
\text{ESPE} &= 1200 \log_2\left(\frac{L_{\text{eff},I}}{L_{\text{eff},II}}\right) \\
&= 1200 \log_2\left(\frac{L_{\text{eff},I}}{L_{\text{eff},I} + \Delta\Delta L}\right)
\end{align}

\subsubsection{Interpretación Física}

ESPE está relacionado con:
\begin{itemize}
    \item Cambios en la corrección de extremos entre octavas
    \item Comportamiento diferente de agujeros en diferentes frecuencias
    \item Efectos de acoplamiento modal
\end{itemize}

\selectlanguage{french}
\subsection{ESPE (Extension Effective du Chemin Sonore)}

% ==================== SECCIÓN COMPLETA: Q-FACTOR ====================

\selectlanguage{spanish}
\subsection{Q-Factor (Factor de Calidad)}

\subsubsection{Definición}

El factor de calidad Q mide la agudeza de una resonancia:

\begin{equation}
Q = \frac{\freqzero}{\Delta f}
\label{eq:q_factor}
\end{equation}

donde $\Delta f$ es el ancho de banda a -3 dB (mitad de potencia).

\subsubsection{Derivación desde la Teoría de Resonadores}

Para un resonador con pérdidas, la admitancia cerca de la resonancia puede modelarse como:

\begin{equation}
\admittance(f) \approx \frac{\admittance_{\max}}{1 + j Q \left(\frac{f}{\freqzero} - \frac{\freqzero}{f}\right)}
\end{equation}

El módulo de la admitancia es:
\begin{equation}
|\admittance(f)| = \frac{\admittance_{\max}}{\sqrt{1 + Q^2 \left(\frac{f}{\freqzero} - \frac{\freqzero}{f}\right)^2}}
\end{equation}

A -3 dB (mitad de potencia), tenemos:
\begin{equation}
|\admittance(f)| = \frac{\admittance_{\max}}{\sqrt{2}}
\end{equation}

Por lo tanto:
\begin{equation}
1 + Q^2 \left(\frac{f}{\freqzero} - \frac{\freqzero}{f}\right)^2 = 2
\end{equation}

Resolviendo para $f$ cerca de $\freqzero$:
\begin{equation}
Q^2 \left(\frac{f - \freqzero}{\freqzero}\right)^2 \approx 1
\end{equation}

\begin{equation}
\frac{f - \freqzero}{\freqzero} \approx \pm \frac{1}{Q}
\end{equation}

El ancho de banda es:
\begin{equation}
\Delta f = 2 \frac{\freqzero}{Q}
\end{equation}

Por lo tanto:
\begin{equation}
Q = \frac{\freqzero}{\Delta f/2} = \frac{2\freqzero}{\Delta f}
\end{equation}

En la práctica, medimos $\Delta f$ como la diferencia entre las frecuencias donde $|\admittance|$ cae a $\admittance_{\max}/\sqrt{2}$.

\subsubsection{Interpretación}

\begin{itemize}
    \item \textbf{Q alto} ($Q > 50$): Resonancia estrecha, color tonal más puro, menos amortiguamiento
    \item \textbf{Q medio} ($10 < Q < 50$): Resonancia moderada, balance entre pureza y riqueza
    \item \textbf{Q bajo} ($Q < 10$): Resonancia ancha, color tonal más rico, más amortiguamiento
\end{itemize}

\selectlanguage{french}
\subsection{Facteur de Qualité Q}

% ==================== SECCIÓN COMPLETA: FRECUENCIAS VS TEMPERAMENTO ====================

\selectlanguage{spanish}
\subsection{Frecuencias de Resonancia vs Temperamento Igual}

\subsubsection{Definición}

Esta métrica compara las frecuencias de resonancia medidas con las frecuencias del temperamento igual (equal temperament):

\begin{equation}
\text{Desviación} = 1200 \log_2\left(\frac{\freqzero}{\freq_{\text{temp}}}\right) \quad \text{[cents]}
\label{eq:deviation_temperament}
\end{equation}

\subsubsection{Derivación del Temperamento Igual}

En el temperamento igual, la frecuencia de una nota se calcula como:

\begin{equation}
\freq_{\text{temp}} = \freq_{\text{La}} \cdot 2^{n/12}
\end{equation}

donde:
\begin{itemize}
    \item $\freq_{\text{La}}$ es la frecuencia del La de referencia (por defecto 415 Hz para barroco)
    \item $n$ es el número de semitonos desde La
\end{itemize}

Para cada nota:
\begin{align}
\text{Re} \quad (n = -7): \quad &\freq_{\text{temp}} = 415 \cdot 2^{-7/12} \approx \SI{311.13}{\hertz} \\
\text{Mi} \quad (n = -5): \quad &\freq_{\text{temp}} = 415 \cdot 2^{-5/12} \approx \SI{349.23}{\hertz} \\
\text{Fa\#} \quad (n = -3): \quad &\freq_{\text{temp}} = 415 \cdot 2^{-3/12} \approx \SI{392.00}{\hertz} \\
\text{La} \quad (n = 0): \quad &\freq_{\text{temp}} = 415 \cdot 2^{0/12} = \SI{415.00}{\hertz}
\end{align}

\subsubsection{Cálculo de Desviación}

La desviación en cents se calcula como:

\begin{align}
\text{Desviación} &= 1200 \log_2\left(\frac{\freqzero}{\freq_{\text{temp}}}\right) \\
&= 1200 \left[\log_2(\freqzero) - \log_2(\freq_{\text{temp}})\right] \\
&= 1200 \left[\log_2(\freqzero) - \log_2(\freq_{\text{La}}) - \frac{n}{12}\right]
\end{align}

\subsubsection{Interpretación}

\begin{itemize}
    \item \textbf{Desviación = 0 cents}: Frecuencia perfectamente temperada
    \item \textbf{Desviación > 0 cents}: Nota más aguda que el temperamento
    \item \textbf{Desviación < 0 cents}: Nota más grave que el temperamento
\end{itemize}

Esta métrica es fundamental para evaluar la afinación del instrumento respecto a estándares musicales.

\selectlanguage{french}
\subsection{Fréquences de Résonance vs Tempérament Égal}

% ==================== SECCIÓN COMPLETA: ALTURA DE PICOS ====================

\selectlanguage{spanish}
\subsection{Altura de Picos de Admitancia}

\subsubsection{Definición}

La altura del pico de admitancia en la frecuencia de resonancia $\freqzero$ es:

\begin{equation}
\admittance_{\max} = \max_f |\admittance(f)|
\label{eq:admittance_peak}
\end{equation}

donde la admitancia es:
\begin{equation}
\admittance(f) = \frac{1}{\impedance(f)}
\end{equation}

\subsubsection{Relación con Facilidad de Emisión}

La altura del pico está relacionada con la facilidad de emisión del sonido. Un pico más alto indica:

\begin{itemize}
    \item Mayor transferencia de energía desde el músico al instrumento
    \item Menor impedancia de entrada
    \item Mayor eficiencia acústica
\end{itemize}

Físicamente, esto se relaciona con:
\begin{equation}
P_{\text{acústica}} = \frac{1}{2} |U|^2 \Re(\impedance) = \frac{1}{2} |p|^2 \Re(\admittance)
\end{equation}

donde $P_{\text{acústica}}$ es la potencia acústica generada.

\selectlanguage{french}
\subsection{Hauteur des Pics d'Admittance}

% ==================== SECCIÓN COMPLETA: RATIO ARMÓNICOS ====================

\selectlanguage{spanish}
\subsection{Ratio de Armónicos Pares/Impares}

\subsubsection{Definición}

El ratio de armónicos pares/impares caracteriza el contenido espectral del sonido:

\begin{equation}
R = \frac{A_{\text{pares}}}{A_{\text{impares}}} = \frac{A_2 + A_4 + A_6 + \ldots}{A_1 + A_3 + A_5 + \ldots}
\label{eq:harmonic_ratio}
\end{equation}

donde $A_n$ es la amplitud del $n$-ésimo armónico.

\subsubsection{Cálculo desde Admitancia}

Las amplitudes se extraen de los picos de admitancia en las frecuencias de antiresonancia:

\begin{align}
A_1 &= |\admittance(\freqzero)| \\
A_2 &= |\admittance(\freqone)| \\
A_3 &= |\admittance(\freqtwo)|
\end{align}

donde $\freqtwo$ es la tercera frecuencia de antiresonancia.

El ratio se calcula como:
\begin{equation}
R = \frac{A_2}{A_1 + A_3}
\end{equation}

\subsubsection{Interpretación Musical}

\begin{itemize}
    \item \textbf{$R > 1$}: Predominio de armónicos pares, sonido más brillante
    \item \textbf{$R < 1$}: Predominio de armónicos impares, sonido más oscuro
    \item \textbf{$R \approx 1$}: Balance entre armónicos, sonido equilibrado
\end{itemize}

\selectlanguage{french}
\subsection{Rapport Harmoniques Pairs/Impairs}

% ==================== SECCIÓN COMPLETA: COHERENCIA DE FASE ====================

\selectlanguage{spanish}
\subsection{Coherencia de Fase}

\subsubsection{Definición}

La coherencia de fase mide la diferencia de fase entre armónicos:

\begin{equation}
\Delta\phi = \arg(\impedance(\freqone)) - \arg(\impedance(\freqzero))
\label{eq:phase_coherence}
\end{equation}

donde $\arg(\impedance)$ es la fase de la impedancia compleja.

\subsubsection{Relación con Claridad del Sonido}

La fase de la impedancia está relacionada con la respuesta temporal del instrumento. Una mayor coherencia de fase indica:

\begin{itemize}
    \item Mayor sincronización entre armónicos
    \item Sonido más claro y definido
    \item Menor distorsión temporal
\end{itemize}

Físicamente, la fase está relacionada con el retardo de grupo:
\begin{equation}
\tau_g = -\frac{d\phi}{d\omega}
\end{equation}

\selectlanguage{french}
\subsection{Cohérence de Phase}

% ==================== SECCIÓN COMPLETA: ESTABILIDAD DE PITCH ====================

\selectlanguage{spanish}
\subsection{Estabilidad de Pitch}

\subsubsection{Definición}

La estabilidad de pitch se basa en la pendiente de la fase cerca de la resonancia:

\begin{equation}
S = \left|\frac{d\phi}{df}\right|_{f=\freqzero}
\label{eq:pitch_stability}
\end{equation}

\subsubsection{Derivación}

Para un resonador, la fase de la impedancia cerca de la resonancia puede aproximarse como:

\begin{equation}
\phi(f) \approx \arctan\left(Q\left(\frac{f}{\freqzero} - \frac{\freqzero}{f}\right)\right)
\end{equation}

La derivada es:
\begin{equation}
\frac{d\phi}{df} = \frac{Q/\freqzero}{1 + Q^2\left(\frac{f}{\freqzero} - \frac{\freqzero}{f}\right)^2}
\end{equation}

Evaluando en $f = \freqzero$:
\begin{equation}
S = \left|\frac{d\phi}{df}\right|_{f=\freqzero} = \frac{Q}{\freqzero}
\end{equation}

\subsubsection{Interpretación}

\begin{itemize}
    \item \textbf{$S$ alto}: Mayor estabilidad, menor sensibilidad a variaciones de presión
    \item \textbf{$S$ bajo}: Menor estabilidad, mayor sensibilidad a variaciones
\end{itemize}

\selectlanguage{french}
\subsection{Stabilité de Hauteur}

% ==================== SECCIÓN COMPLETA: CUT-OFF FREQUENCY ====================

\selectlanguage{spanish}
\subsection{Frecuencia de Corte (Cut-off Frequency)}

\subsubsection{Definición}

La frecuencia de corte es la frecuencia más alta donde la admitancia normalizada cae por debajo de un umbral:

\begin{equation}
\freq_{\text{cut}} = \max\left\{f : \frac{|\admittance(f)|}{|\admittance_{\max}|} \geq \theta\right\}
\label{eq:cutoff_freq}
\end{equation}

donde $\theta$ es el umbral (típicamente 0.1 = 10\%).

\subsubsection{Interpretación Física}

La frecuencia de corte está relacionada con:

\begin{itemize}
    \item Límite superior del rango útil del instrumento
    \item Capacidad de generar armónicos superiores
    \item Eficiencia de radiación a altas frecuencias
\end{itemize}

Físicamente, está relacionada con el diámetro del conducto:
\begin{equation}
\freq_{\text{cut}} \propto \frac{c_0}{d}
\end{equation}

donde $d$ es el diámetro característico del conducto.

\selectlanguage{french}
\subsection{Fréquence de Coupure}

% ==================== NOTA FINAL ====================
% Estas secciones deben integrarse en el documento principal
% en el orden apropiado dentro de la Sección 4 (Mediciones Acústicas)

